\section{Problem Formulation}
\label{sec:problem}

We consider multistage constant-current (MCC) charging from an initial SOC to
a prescribed target. Each protocol is evaluated by its charging time, peak
temperature rise, and charging-induced degradation. Figure~\ref{fig:charging_problem}
summarizes the decision variables, objectives, and constraints.

\begin{figure}[!t]
\centering
\includegraphics[width=\columnwidth]{fig1.pdf}
\caption{Fast-charging multiobjective design problem.}
\label{fig:charging_problem}
\end{figure}


\subsection{Multistage Constant-Current Protocol}

A $K$-stage protocol applies a current $I_k$ over an SOC increment $\Delta s_k$
at stage $k$. The initial and target SOCs are $s_0$ and $s_f$, respectively.
The independent decision variables are
\begin{equation}
\theta=[I_1,\ldots,I_K,\Delta s_1,\ldots,\Delta s_{K-1}]^\top
\label{eq:protocol}
\end{equation}
and the final SOC increment is
\begin{equation}
\Delta s_K=s_f-s_0-\sum_{k=1}^{K-1}\Delta s_k
\label{eq:terminal_soc}
\end{equation}
Each current and SOC increment is restricted to its prescribed interval. The
last increment must also exceed a positive lower bound so that every stage has
nonzero duration.
% AUTHOR_CHECK P1: Supply the actual current/SOC bounds, terminal lower bound,
% initial and target SOCs, and effective capacities from the experiment configs.

Following~\cite{nref7,nref8}, we optimize SOC increments rather than stage
durations. The current and charge delivered in each stage determine its
duration, giving the total charging time
\begin{equation}
t_c(\theta)=3600Q_{\mathrm{eff}}\sum_{k=1}^{K}\frac{\Delta s_k}{I_k}
\label{eq:charging_time}
\end{equation}
where $Q_{\mathrm{eff}}$ is the effective cell capacity in Ah and current is
measured in A. The experiments use $K=3$, giving five independent decision
variables.

\subsection{Constrained Multiobjective Charging Design}

The three objectives are minimized jointly:
\begin{equation}
\min_{\theta}\quad
\mathbf f(\theta)=[t_c(\theta),\,T_p(\theta),\,Q_s(\theta)]^\top
\label{eq:cmop}
\end{equation}
Here, $T_p$ is the peak temperature rise and $Q_s$ is the aging indicator
defined in Section~\ref{sec:model_section}. In addition to the bounds on the
decision variables, the charging trajectory must satisfy
\begin{equation}
U(t;\theta)\le U_{\max},\quad
T(t;\theta)\le T_{\max},\quad
0\le t\le t_c(\theta)
\label{eq:operating_constraints}
\end{equation}
where $U_{\max}$ and $T_{\max}$ are the voltage and temperature limits. A
protocol is infeasible if it violates either limit or fails to reach $s_f$;
such protocols are excluded from the Pareto archive.
% AUTHOR_CHECK P2: Exclusion from the Pareto archive does not specify how
% infeasible/failed simulations are represented in the GP or acquisition.

Higher charging currents reduce charging time but generally increase the
thermal and aging costs. We therefore seek a set of feasible nondominated
protocols, allowing the charging-time, temperature, and degradation trade-off
to be chosen for the operating requirement.
